\documentclass[11pt,letterpaper]{article}
\PassOptionsToPackage{hyphens}{url}
\usepackage[margin=.83in,top=.8in,bottom=.8in,headheight=15pt]{geometry}
\usepackage[T1]{fontenc}
\usepackage{lmodern,amsmath,amssymb,mathtools,bm,booktabs,array,graphicx,tikz}
\usetikzlibrary{arrows.meta,positioning,calc,decorations.pathreplacing,patterns}
\definecolor{ink}{HTML}{17344D}\definecolor{accent}{HTML}{087E8B}\definecolor{pale}{HTML}{EDF5F5}\definecolor{gold}{HTML}{AD6B19}
\providecommand{\dd}{\,\mathrm{d}}\providecommand{\R}{\mathbb{R}}
\providecommand{\E}{\mathbb{E}}\providecommand{\Var}{\operatorname{Var}}
\providecommand{\Cov}{\operatorname{Cov}}\providecommand{\diag}{\operatorname{diag}}
\providecommand{\norm}[1]{\left\lVert#1\right\rVert}
\providecommand{\abs}[1]{\left|#1\right|}
\usepackage{microtype,hyperref,fancyhdr,enumitem,amsthm,sectsty,needspace}
\allsectionsfont{\raggedright}
\hypersetup{colorlinks=true,linkcolor=ink,urlcolor=accent,pdfauthor={Zan Ahmad},pdftitle={Coordinates, sampled motion, and Python expressions / Review}}
\setlength{\parindent}{0pt}\setlength{\parskip}{6pt}\linespread{1.07}
\setlist{itemsep=3pt,topsep=4pt,leftmargin=*}
\emergencystretch=2em\clubpenalty=10000\widowpenalty=10000
\pagestyle{fancy}\fancyhf{}\fancyhead[L]{\small Scientific Computing in Python / 2027}\fancyhead[R]{\small Lecture 1}
\fancyfoot[L]{\footnotesize Zan Ahmad}\fancyfoot[R]{\footnotesize\thepage}\renewcommand{\headrulewidth}{.3pt}
\newtheorem{proposition}{Proposition}\newtheorem{theorem}{Theorem}\newtheorem{lemma}{Lemma}
\theoremstyle{definition}\newtheorem{definition}{Definition}\newtheorem{example}{Example}
\numberwithin{equation}{section}
\begin{document}
\section*{Coordinates, sampled motion, and Python expressions --- review sheet}
Use this after attempting the worked examples and practice. Explain each claim in words before calculating.

\section*{Key ideas}
\begin{itemize}
\item Choose the quantity before choosing a command.
\item Displacement subtracts the starting coordinate.
\item Average velocity does not reveal unobserved path length.
\item Assignment stores a value at execution time.
\item Origin shifts and unit conversions give predictable checks.
\item A clean run and a bounded explanation belong together.
\end{itemize}
\section*{Common errors and repairs}
\textbf{Divide final position by elapsed time.} Subtract the starting position first.

\textbf{Call displacement/time the true average path speed.} It is signed interval-average velocity; path length may be unknown.

\textbf{Change a unit label alone.} Convert the numerical values too.

\section*{Retrieve without your notes}
\begin{itemize}
\item A moves from 5 to -1 mm in 3 s. What is its average velocity?
\item What does an origin shift change?
\item What does a clean kernel run establish?
\end{itemize}
\section*{Readiness check}
\begin{itemize}
\item Separate final position, displacement and interval-average velocity in a stated question.
\item Use numerical values, text labels, assignment and list indices to calculate a quantity with units.
\item Predict how changing a coordinate origin or length unit changes a result.
\item Distinguish observed points from unobserved motion between samples.
\item Restart and rerun a short notebook and explain its scientific limits.
\end{itemize}
\end{document}
